\documentclass[11pt,a4paper]{article}
\usepackage[margin=24mm]{geometry}
\usepackage[T1]{fontenc}
\usepackage{lmodern,amsmath,amssymb,booktabs,tabularx,hyperref}
\hypersetup{colorlinks=true,urlcolor=blue,citecolor=blue}
\title{Device-Aware Low-Order Floating-Point Coupling\\for MRAM-Based p-Bit Ising Machines}
\author{Du Yanzhang\\National University of Singapore\\Supervisor: Prof. Kelvin Fong}
\date{Research scope and M0 protocol --- 6 November 2026}
\begin{document}
\maketitle
\begin{abstract}
This proposal develops low-order floating-point coupling formats (FP4, FP6, FP8) with configurable exponent scale granularity for nonideal MRAM p-bit systems, and evaluates their quality-cost tradeoff on independent BPSK MIMO detection. The design target is a device-aware low-order floating-point coupling architecture, supported by a configuration rule, bit-true macro and synthesis. Optimized scaled integer and block-floating formats are mandatory cost-matched controls; no FP superiority is assumed. All previous computational results are excluded from the new evidence base.
\end{abstract}
\section{Locked research question}
\emph{How should coupling precision and scale granularity be chosen for nonideal MRAM p-bit systems to preserve useful optimization performance at lower implementation cost?}
\section{Reference mathematical model}
An ideal p-bit network uses binary variables with input-controlled stochastic response \cite{camsari2017}. For symmetric $J_{ij}=J_{ji}$ and $J_{ii}=0$:
\begin{align}
E(\mathbf{s}) &= C-\sum_{i<j}J_{ij}s_is_j-\sum_i h_i s_i,\\
H_i &= h_i+\sum_{j\ne i}J_{ij}s_j,\\
P(s_i=+1\mid\mathbf{s}_{-i}) &= \left[1+\exp(-2\beta H_i)\right]^{-1}.
\end{align}
This asynchronous Gibbs model is mathematical reference software, not a validated MTJ device or physical circuit.
\section{Locked hypotheses}
\begin{description}
\item[H1: Device-dependent useful precision.] Does minimum sufficient coupling precision change with the type and severity of p-bit nonideality? Define an absolute task-quality floor.
\item[H2: Scale granularity.] When do per-element exponents justify their storage and arithmetic cost compared with optimized global, row, tile, or block scales? Preserve symmetric effective couplings.
\item[H3: Size dependence.] Across BPSK MIMO channel dimensions, antenna ratios, SNR and conditioning, determine useful precision requirements without assuming a trend.
\end{description}
\section{Proposed Coupling Array and Relationship to CIM}
The FYP retains an array-based hardware target: a stochastic $p$-bit Ising machine with a low-order floating-point coupling array, not the earlier conventional MRAM-CIM current-summing/ADC array. For $N=N_t$ BPSK MIMO spins, the logical architecture has a symmetric $N\times N$ coupling matrix $J$, $N$ biases $h_i$, and a feedback spin state vector. The coupling hardware computes a selected local field
\begin{equation}
H_i=h_i+\sum_{j\ne i}J_{ij}s_j,
\end{equation}
using stored low-order FP4/FP6/FP8 coefficients, sign selection, exponent/scale alignment and bounded-width accumulation. A stochastic MRAM $p$-bit updates its state according to its actual probability-controlling input. The ideal sigmoid of the mathematical reference is not a foundry-validated magnetic-device response.

The M2--M3 hardware milestone remains a \emph{bit-true 64--128-spin coupling tile} and synthesis-backed area/timing estimate. This tile includes coupling/bias storage assumptions, format decoding, accumulation, update scheduling, p-bit-interface abstraction and accounting for interconnect, bandwidth and scale metadata. It does \emph{not} imply a fabricated 64--128-$p$-bit MTJ array or rebuilding the former 256$\times$256 2T1MTJ-CIM circuit.

Three physical integrations must be distinguished: a digital SRAM/register-driven coupling processor is a $p$-bit accelerator; MRAM-backed coefficient storage with adjacent arithmetic is near-memory probabilistic processing (its CIM label depends on integration); and a true in-array weighted-sum circuit can qualify as CIM. Signed FP exponents and local-field accumulation require explicit circuitry in each case---ordinary low-bit FP coefficients do not map directly onto analog MTJ conductances. Crossbar--SMTJ probabilistic machines provide relevant architectural precedent but \emph{not} low-order-FP implementation proof \cite{hybrid2026,vcma2024}.

Thus the novelty remains \emph{low-order floating-point coupling under nonideal MRAM $p$-bit behavior}, with BPSK MIMO as the workload and a defensible array/compute boundary established before claiming CIM.
\section{Fair comparisons}
Compare high-precision reference, optimized scaled integers, shared-scale block floating point, and precisely defined per-element low-bit float. Count scales, metadata, exponent alignment, accumulation, bandwidth and replica overhead. Assess both $J$ and $h$ quantization; isolated $J$-only analysis is explicitly labeled. Keep local-field updates and ideal replica exchanges consistent with the quantized Hamiltonian. Score all outputs against the unquantized BPSK MIMO ML objective; report BER against transmitted symbols separately. The core design targets FP4 (E2M1), FP6 (E3M2), FP8 (E4M3) coupling and exponent handling. Historical TSP and graph benchmarks are removed from formal application evaluation.
\section{BPSK MIMO benchmark}
For complex-baseband BPSK, $\mathbf{y}=\mathbf{H}\mathbf{x}+\mathbf{n}$, $x_i\in\{-1,+1\}$, and maximum likelihood detection minimizes $R(\mathbf{s})=\|\mathbf{y}-\mathbf{H}\mathbf{s}\|_2^2$. With $G=\Re(\mathbf{H}^{\mathrm H}\mathbf{H})$ and $b=\Re(\mathbf{H}^{\mathrm H}\mathbf{y})$, set $J_{ij}=-2G_{ij}$ for $i\ne j$, $h_i=2b_i$, and $C=\|\mathbf{y}\|_2^2+\sum_iG_{ii}$, so the Ising energy equals the original ML residual. Small cases are exhaustively checked; larger systems have no assumed exact oracle. Prior work already studies finite-precision MIMO Ising couplings \cite{singh2022}; our novelty is the device-aware low-order FP coupling macro.
\section{Milestones M0--M5}
\begin{center}\small
\begin{tabularx}{\linewidth}{@{}lX@{}}\toprule
Stage & Completion requirement\\\midrule
M0 & Exact MIMO-to-Ising reference, seeded BPSK channels, BER and nonideal p-bit response-model provenance.\\
M1 & FP4/FP6/FP8 versus tuned scaled integers and BFP under device response, paired uncertainty.\\
M2 & Device-aware selection rule and coupling-macro specification with full analytical cost.\\
M3 & Bit-accurate 64--128-spin tile, synthesis-backed area and timing comparison.\\
M4 & Held-out BPSK MIMO instances, device populations and supported scaling limits.\\
M5 & Final architecture, configuration rules and reproducible report.\\\bottomrule
\end{tabularx}
\end{center}
\section{Evaluation and evidence}
Primary: time-to-target on the original ML objective, counting all updates, exchanges and selection. Report bit-error rate against transmitted BPSK symbols versus SNR separately. Secondary: original-task objective gap, target-hit probability, validity, ideal small-$N$ distribution check and verified implementation cost. Software update counts do not establish hardware latency or energy. Independent instances are the primary statistical unit; solver seeds are nested. An unvalidated behavioral sigmoid does not constitute physical MRAM evidence. All older CIM, sensing-margin, TSP pilot and assistant-generated nonideality results are excluded.
\section{Preliminary low-order floating-point coupling screen (E1/E2)}
The first fresh NUS BPSK MIMO software screen used a complex $8\times8$ channel,
$12$ independent channel/noise frames per SNR ($96$ transmitted bits),
SNR $0$ and $5$~dB, and $1000$ attempted asynchronous spin updates per
format/response/frame. All 24 then-current unit tests passed.
The screening combined preliminary custom low-order coupling codebooks and
\emph{illustrative}, not experimentally validated, p-bit response variants.
Table~\ref{tab:m1pilot} reports independent-frame exact-ML solution hits
under the ideal response; the source is the user-executed NUS M1 software
log dated 6 November 2026, not magnetic-device or synthesis evidence.
\begin{table}[htbp]
\centering\small
\begin{tabular}{@{}lrrr@{}}\toprule
Representation & 0 dB ML hits & 5 dB ML hits & 0 dB bit errors\\\midrule
Unquantized reference & 12/12 & 12/12 & 7/96\\
FP4 E2M1 & 10/12 & 10/12 & 9/96\\
INT4 scaled & 11/12 & 12/12 & 9/96\\
BFP4 shared-scale & 10/12 & 10/12 & 8/96\\
FP6 E3M2 & 12/12 & 12/12 & 7/96\\
INT6 scaled & 12/12 & 12/12 & 7/96\\
FP8 E4M3 & 12/12 & 12/12 & 7/96\\
INT8 scaled & 12/12 & 12/12 & 7/96\\\bottomrule
\end{tabular}
\caption{Early paired M1 screen (12 frames per SNR). ML-state hit
and transmitted-bit BER are separate metrics; sample counts are insufficient
for inferential hardware or precision claims.}
\label{tab:m1pilot}
\end{table}
Across all formats, the illustrative gain-0.80 and fixed per-device
offset-$\sigma=0.15$ conditions gave the same aggregated BER/ML-hit
outcomes as ideal; the contrast-0.90 condition altered only the BFP4
0~dB outcomes. In that case, BFP4 ML hits decreased from $10/12$ to
$9/12$, while observed transmitted-bit errors decreased from $8/96$
to $7/96$, illustrating why BER alone is an unreliable optimization
proxy in a small pilot. This evidence does \emph{not} establish the
hypothesized format--nonideality interaction.

An important reference erratum is that the pilot's
\texttt{fp32}-labeled path retained Python double-precision values while
counting 32 nominal bits per coefficient. It should be described as an
\emph{unquantized software reference}, not a measured binary32 storage
implementation. Subsequent code separates genuine binary32 coefficient
rounding from a 64-bit unquantized reference. Local-field accumulation
still requires bit-true modeling. Further limitations include per-frame
coefficient-MSE scale optimization, non-equal scale granularity across
FP/INT and BFP, tiny frame samples, full-budget best-of-sample selection,
and lack of a physically characterized stochastic MRAM response.
These results motivate paired final-state/first-hit diagnostics and
a deployable scale/accumulator study; they do not motivate changing
the locked BPSK MIMO benchmark or low-order-FP hardware direction.

\subsection{Paired response sensitivity and exact-ground diagnostic}
The same NUS pilot was subsequently analyzed in paired comparisons against
the ideal response (12 original frames per SNR and representation).
The illustrative gain-$0.80$ and offset-$\sigma=0.15$ variants changed
\emph{no final selected states} among the tested 10 format labels and both
SNRs, though some first-best-recorded checkpoints changed.
The contrast-$0.90$ response changed only one selected state:
BFP4 at $0$~dB (1 of 12 paired frames), and reduced that group's
original-ML hit count by one. These final-state comparisons are highly
dependent across formats, response settings and SNRs, so they are not
independent experimental replicates. A different first-best checkpoint
is \emph{not} a measured hardware time-to-target.

At the same four-bit nominal storage, paired original-ML hit comparisons
favored INT4 in 2 frames versus FP4 in 1 frame at $0$~dB, and
in 2 frames versus 0 at $5$~dB; the total bit-error counts were
identical. No meaningful FP-versus-INT superiority is established.

\subsection{Exact quantized-ground audit: representation versus sampling}
The exact audit was subsequently completed on the original 6 November
NUS MIMO pilot, independently enumerating all $2^8=256$ states
for each quantized Hamiltonian, without rerunning stochastic sampling.
The recorded selected state minimized its \emph{own quantized}
Hamiltonian in every tested frame/format/response combination
($12/12$ per row, within the numerical tolerance used by the audit).
Therefore the observed original-ML failures in this pilot cannot be
attributed to failure to reach a \emph{quantized} minimum.

\begin{table}[htbp]
\centering\small
\begin{tabular}{@{}lrr@{}}\toprule
Representation & 0 dB preserved & 5 dB preserved\\\midrule
FP4 E2M1 & 10/12 & 10/12\\
INT4 scaled & 11/12 & 12/12\\
BFP4 shared-scale & 11/12 & 11/12\\
FP6 / FP8 & 12/12 & 12/12\\
INT6 / INT8 & 12/12 & 12/12\\
BFP6 / BFP8 & 12/12 & 12/12\\\bottomrule
\end{tabular}
\caption{Original unquantized ML state remains among minimizers of
the quantized Hamiltonian. Values use the \emph{ideal} response;
the quantized coefficient landscape is unchanged by the illustrative
p-bit response model. Each SNR has only 12 frames and the two SNRs
share channel seeds.}
\label{tab:m1ground}
\end{table}

For FP4 the original ML optimum was absent from the quantized
ground-state set in $2/12$ frames at each SNR. INT4 lost it
in $1/12$ frames at $0$~dB and $0/12$ at $5$~dB.
For BFP6 at $5$~dB, the original ML state was retained
in $12/12$ quantized ground-state sets but the selected state
matched original ML in only $11/12$; the same phenomenon
occurs for BFP4. This is consistent with \emph{quantization-induced
degeneracy}: multiple equal-minimum quantized states can have unequal
original ML residuals, with stochastic search selecting one of them.
A subsequent read-only NUS CSV audit confirmed the ground-state
multiplicity pattern: FP4 and INT4 had no degenerate cases at either
SNR; BFP4 had two such frames at 0~dB (maximum two minima)
and one at 5~dB (maximum three minima); BFP6 had one case at
5~dB (maximum two minima), and none at 0~dB.
More specifically, the original ML state remained a quantized minimizer
but was not selected in BFP4 0~dB frame 10, BFP4 5~dB frame 8,
and BFP6 5~dB frame 10. These are frame identifiers from the
original NUS experiment, not independent replicated trials.
Here equality means energies tied within the exact audit's
tolerance ($10^{-8}$); a bit-accurate accumulator check will be
needed to determine whether the corresponding physical machine
would realize the same degeneracy.

H1 is \emph{not validated}: gain-$0.80$ and offset-$0.15$
illustrative responses did not change final selections, and
contrast-$0.90$ changed one BFP4 selected state at $0$~dB.
While some best-recorded checkpoint indices changed, these
are not calibrated time-to-target. H2 is not isolated either:
FP/INT used an optimistic per-frame global scale, whereas
BFP used tile-pair power-of-two scales. A necessary next study
must control scale granularity, energy degeneracy and finite-budget
search, while separately establishing a physically sourced
stochastic MRAM p-bit response. No FP superiority, CIM implementation,
hardware energy or novelty claim follows from this small software pilot.
\subsection{Planned finite-only floating-point and $J$/$h$ sensitivity gate}
To separate representation-induced objective changes from stochastic
sampling, the next experiment considers the finite-only FP4 E2M1 and
FP6 E3M2 \emph{element formats} standardized by OCP \cite{ocpmx2023}.
Finite-only E2M1 supports magnitudes $0,0.5,1,1.5,2,3,4,6$;
historical FP4 (kept for reproducibility) only reaches magnitude $3$.
The format itself is established prior art, and its adoption is not
equivalent to a complete MX block-scaled or bit-exact circuit.

New FP and integer coupling scales are calibrated only on disjoint
training BPSK MIMO frames, then frozen for held-out frames, with
identical global/tile grouping. J-only and h-only quantization are
\emph{unquantized-opposite oracle diagnostics}, while mixed FP4-J/FP6-h
is compared against equally mixed INT4-J/INT6-h with all coefficient
and scale storage counted. Upgrading the $N$ bias coefficients from
four to six bits adds $2N$ nominal coefficient bits (256 at $N=128$),
before decoding, accumulation, control logic and other overhead.

These comparisons were later executed as fresh NUS M1B/M1C
held-out exact-enumeration software experiments, summarized below.
H1 remains open: finite-budget original-task performance must be studied
under physically sourced stochastic MRAM p-bit characteristics, and
M2/M3 still require bit-true arithmetic and RTL synthesis.
\subsection{M1B: first held-out frozen-scale results (E0/E1)}
A separate NUS user-run experiment applied 24 calibration and
24 disjoint held-out channel frames per SNR (0 and 5~dB) for
$8\times8$ BPSK MIMO. A group scale was fitted using calibration
coefficients only, then frozen for the held-out Hamiltonians.
All outcomes in Table~\ref{tab:m1b-heldout} are \emph{exact enumeration}
of quantized-Hamiltonian minima, not measured p-bit detection BER,
nonideal device response or hardware-area estimates. The software
reported 36 then-current tests passing, indicating this was the
older M1B implementation; in particular finite-only \texttt{fp4\_fn}
and \texttt{fp6\_fn} and the subsequent J/h ablation were \emph{not}
run in this dataset.

\begin{table}[htbp]
\centering\small
\begin{tabular}{@{}lrrrr@{}}\toprule
Format & Global ML 0/5~dB & Tile ML 0/5~dB & Global bits & Tile bits\\\midrule
Historical FP4 & 18/24, 21/24 & 18/24, 21/24 & 208 & 272\\
INT4 & 18/24, 21/24 & 18/24, 22/24 & 208 & 272\\
BFP4 & 19/24, 21/24 & 19/24, 21/24 & 160 & 176\\
Historical FP6 & 21/24, 24/24 & 21/24, 24/24 & 280 & 344\\
INT6 & 21/24, 24/24 & 23/24, 22/24 & 280 & 344\\
BFP6 & 23/24, 24/24 & 23/24, 24/24 & 232 & 248\\
Historical FP8 & 24/24, 24/24 & 24/24, 24/24 & 352 & 416\\
INT8 & 24/24, 24/24 & 24/24, 24/24 & 352 & 416\\
BFP8 & 24/24, 24/24 & 24/24, 24/24 & 304 & 320\\\bottomrule
\end{tabular}
\caption{First frozen-CAL/held-out exact original-ML preservation
(retained cases out of 24 per SNR), with nominal unique $J$/$h$
coefficient plus assumed scale-metadata bits. The 0 and 5~dB
cases use paired channel seeds, not 48 independent frames; the
scale representation is not bit-true nor hardware synthesized.}
\label{tab:m1b-heldout}
\end{table}

Historical FP4 and INT4 were equal under global scaling; INT4
retained one extra 5~dB case under tile scaling. BFP6 preserved
23/24 original ML states at 0~dB compared with 21/24 for
historical FP6 and INT6 globally, using fewer nominal bits.
This is a small sample and uses different scale encodings
(BFP uses power-of-two scales with eight-bit exponent metadata,
FP/INT assume 32-bit real scales); it is not an equal-area or
statistically established design advantage. Eight-bit formats
reached an apparent ceiling at 24/24 in this task.
Nearly identical J-MSE across SNR for each format/group follows
from reusing the same channel realizations and $J$ independent
of received noise. Local field biases $h$ and the ML landscape
can still differ across SNR. Subsequent M1B/M1C codebook
and bias/coupling ablations are described in the next subsection.
\subsection{M1B/M1C: held-out finite-only FP and $J$/$h$ ablation}
The next NUS software run retained the same complex $8\times8$
BPSK MIMO benchmark with 24 independent calibration and 24
disjoint held-out channel seeds, tested at 0 and 5~dB SNR.
The two SNRs are paired on the same channel seeds. Coupling/bias
scales were fitted by training-set coefficient MSE and \emph{frozen}
for held-out testing. The 0~dB and 5~dB original-ML preservation
counts in Table~\ref{tab:m1c-mixed} are \emph{exact representational
ground-state} results and \emph{not} finite-budget p-bit sampling
success, measured stochastic MTJ data, or statistical evidence
of FP superiority.

\begin{table}[htbp]
\centering\small
\begin{tabular}{@{}lrrr@{}}\toprule
Global-scale format & 0~dB & 5~dB & Nominal bits\\\midrule
Historical FP4, $J$ and $h$ & 18/24 & 21/24 & 208\\
Finite-only FP4, $J$ and $h$ & 19/24 & 22/24 & 208\\
INT4, $J$ and $h$ & 18/24 & 21/24 & 208\\
\textbf{Finite-only FP4-$J$ / FP6-$h$} & \textbf{23/24} & \textbf{24/24} & 224\\
INT4-$J$ / INT6-$h$ & 21/24 & 24/24 & 224\\
Finite-only FP6, $J$ and $h$ & 21/24 & 24/24 & 280\\
BFP6, $J$ and $h$ & 23/24 & 24/24 & 232\\\bottomrule
\end{tabular}
\caption{Held-out original ML solution preserved within the
quantized Hamiltonian's ground-state set. Nominal bits count
unique undirected $J$, $h$, and assumed scale metadata; they do
not measure local-field arithmetic, circuitry, area, or energy.
Twenty-four cases per SNR are insufficient for an FP
superiority claim.}
\label{tab:m1c-mixed}
\end{table}

The $J$-only/$h$-only oracles help localize the damage:
finite-only FP4 at 0~dB retained 23/24 when only $J$ was
quantized (exact $h$), versus 19/24 when only $h$ was quantized
(exact $J$), and 19/24 when both were quantized. Matched INT4
gave 21/24, 18/24, and 18/24 respectively. These are
\emph{unquantized-opposite diagnostics}, not practical low-bit
hardware designs. Their interaction need not be additive.
At 5~dB, finite-only FP4 $J$-only preserved 24/24,
$h$-only 23/24, and both 22/24.

Representing $J$ in finite-only FP4 and $h$ in finite-only
FP6 increased global-scale original-ML preservation to
23/24 at 0~dB and 24/24 at 5~dB with 224 nominal bits.
At the same nominal bits, mixed INT4-$J$/INT6-$h$ retained
21/24 and 24/24; the two-case difference at 0~dB
\emph{does not} establish a statistically robust FP advantage.
Under tile grouping both mixed configurations retained 22/24
and 24/24 at 0/5~dB, respectively, at 288 nominal bits.
BFP6 globally retained 23/24 and 24/24 with 232 nominal bits,
a serious fixed-point/power-of-two-scale baseline for hardware.

Finite-only E2M1 made a small global improvement over the
historical FP4 codebook (19 versus 18 preserved cases
at 0~dB, and 22 versus 21 at 5~dB), but \emph{did not}
consistently improve under tile scaling. An element format
with more representable values does not by itself
guarantee detection quality after scale fitting.
The mixed-FP result is a candidate \emph{architectural
hypothesis}, not device-aware novelty: all of these results use
software-real-valued scales, exact enumeration, and no
physical stochastic MRAM response or synthesized
exponent-alignment/accumulator cost.

To avoid open-ended format sweeps, further implementation
work will focus on three candidate configurations:
finite-only FP4-$J$/FP6-$h$, mixed INT4-$J$/INT6-$h$,
and BFP6. The next decisive evidence is a physically
sourced p-bit response, a fixed finite-budget original
MIMO-quality criterion, and bit-true RTL/synthesis
cost for the proposed 64--128-spin coupling tile.

\subsection{Required physical MRAM p-bit response gate}
Unlike a high-retention static eMRAM storage MTJ, a stochastic
$p$-bit requires a fluctuation mechanism with a controllable
state probability and a characterized update time. Recent
hybrid memristor--stochastic-MTJ experiments measured a
voltage-dependent sigmoidal switching probability, strong
pulse-width sensitivity in slower devices, and variability
among SMTJs \cite{hybrid2026}. Independent work reported a
CMOS-integrated superparamagnetic-MTJ $p$-bit at 130~nm
\cite{cacoilo2026} and an on-chip stochastic-MTJ/transistor
$p$-bit core with material- and circuit-dependent response
nonidealities \cite{daniel2024}. These studies show plausible
response mechanisms, but \emph{do not validate} our
available stable-storage MRAM model as a stochastic MTJ or
establish process portability to the proposed coupling tile.

Before H1 can support a physical novelty claim, we must
identify a permitted stochastic-device response model
or measured voltage/current-to-probability data, its
pulse-width/time correlation and device variation, and
the circuit mapping from accumulated $H_i$ to the
$p$-bit input. If external publications provide only
cross-technology response curves, results must be
explicitly labeled literature-informed behavioral modeling.
The decisive experiment then compares fixed finite
spin-update budgets and original-task MIMO quality for
the selected mixed FP, matched mixed INT, and BFP6
designs with uncertainty and bit-true synthesized cost.
No more unconstrained codebook sweeps are justified
by the existing 24-frame result.

\subsection{H2: held-out 16-QAM encoding-aware scale experiment}
A fresh NUS software experiment verified the complex 16-QAM
spin mapping against all $2^{16}$ original-channel residuals
per 4-by-4 complex MIMO frame, including intra-symbol
significance-bit couplings and Gray data labeling.
Using 100 calibration frames, 300 disjoint held-out frames
at 12~dB, and frozen symmetric global, tile, or
significance-pair grouping, both $J$ and $h$ were quantized.
Exact enumeration gave the following original-ML-state
retention counts: global mixed FP4-$J$/FP6-$h$ 48/300,
matched INT4-$J$/INT6-$h$ 39/300, BFP6 109/300;
significance-pair mixed FP 57/300, mixed INT 41/300,
BFP6 104/300. Global nominal coefficient-plus-assumed
scale storage was 640, 640, and 832 bits, respectively.
Mixed FP showed a numerical preference over matched
INT on these paired frames, but paired uncertainty
has not been computed and no FP advantage is
statistically established. BFP6 had much better
represented ground-state preservation yet 107/300
global frames contained multiple quantized minima,
so the actual stochastic p-bit output-selection
distribution is a critical open question. Uniform
selection among quantized minima gave Gray BER
0.2049 for mixed FP, 0.2240 for matched INT,
and 0.1796 for BFP6 globally; original full-precision
ML Gray BER was 0.1210. These tie-average figures
are mathematical diagnostics, not sampled p-bit BER.

This H2 result (local NUS directory
\texttt{h2\_16qam\_20261106T040315Z\_4108213})
does not include stochastic device characteristics
or synthesized circuit cost. The next H1 experiment
uses the same frozen calibration/representation
implementation with an exact one-p-bit-at-a-time
Markov chain on eight spins, separately recording
probability of visiting an original-quality target
and probability of \emph{returning} one at a
fixed update budget. Synthetic contrast/mismatch
responses are E2 illustrations pending device sourcing.

\subsection{H1: exact eight-spin Markov chain, first NUS result}
An exact random-scan p-bit chain was evaluated on the
$2\times2$ complex 16-QAM (eight-spin) benchmark at 12~dB,
using 20 held-out channels and two illustrative device
populations, four previously frozen coupling representations,
and $\beta=1$. The E2 response variants were an ideal sigmoid,
contrast-$0.90$ response, and heterogeneous static gain/offset.
The exact transition dynamics model verified the ideal Gibbs
stationary distribution; all recorded aggregate cases converged.
At 500 attempted updates, mixed FP4-$J$/FP6-$h$ reached
an original-ML residual-gap target ($\delta=0.05$) at least
once with probability $0.9713$, but its probability of
\emph{being} in that target at the fixed observation time
was only $0.0507$; expected Gray BER was $0.3834$.
Matched mixed INT achieved $0.9709$, $0.0521$, and
$0.3825$; BFP6 achieved $0.9752$, $0.0522$, and $0.3811$;
uniform FP6 achieved $0.9761$, $0.0526$, and $0.3818$.
These results show that visiting a good state and returning
one at a fixed time are different objectives. Neither the
higher precision benefit nor a device-dependent useful
precision rule was established in this small test.

Importantly, $\beta=1$ was an arbitrary illustrative
operating point. With the generator's complex AWGN variance
$\sigma_n^2=(N_t/N_r)10^{-\mathrm{SNR}/10}$, a Gibbs sampler
matching the AWGN posterior $\propto\exp[-R(s)/\sigma_n^2]$
at 12~dB in square MIMO requires
$\beta=1/\sigma_n^2\approx15.85$. Thus, poor fixed-time
quality at $\beta=1$ should not be interpreted as an
intrinsic precision plateau. A single predeclared
posterior-consistent operating-point comparison is needed;
it remains a software-model diagnostic, not a stochastic
MTJ voltage-response validation. The two device populations
do not create 40 independent channel frames, and all
nonideal responses remain synthetic E2.
\section{Project schedule}
Fourteen active working weeks and two weeks buffer, reduced capacity around exams. Aim for M1/M2 for CA3. FPGA deployment is optional. A well-supported integer-favoring configuration rule can be a valid outcome.
\begin{thebibliography}{9}
\bibitem{camsari2017} K. Y. Camsari, R. Faria, B. M. Sutton and S. Datta, ``Stochastic p-Bits for Invertible Logic,'' \emph{Physical Review X} 7, 031014 (2017), DOI: 10.1103/PhysRevX.7.031014.
\bibitem{singh2022} A. K. Singh et al., ``Ising Machines' Dynamics and Regularization for Near-Optimal MIMO Detection,'' \emph{IEEE Transactions on Wireless Communications} 21, 11080--11094 (2022), DOI 10.1109/TWC.2022.3189604.
\bibitem{hybrid2026} ``Intrinsic annealing in a hybrid memristor-magnetic tunnel junction Ising machine,'' \emph{Nature Communications} (2026), DOI: 10.1038/s41467-026-71844-8.
\bibitem{vcma2024} ``In-Memory Neural Stochastic Differential Equations with Probabilistic Differential Pair Achieved by In-Situ P-Bit Using CMOS Integrated Voltage-Controlled Magnetic Tunnel Junctions,'' \emph{IEDM} (2024), DOI: 10.1109/IEDM50854.2024.10873318.
\bibitem{ocpmx2023} Open Compute Project, ``OCP Microscaling Formats (MX) Specification,'' Version 1.0, Sep. 2023, Secs. 5.3.2--5.3.3, \url{https://www.opencompute.org/documents/ocp-microscaling-formats-mx-v1-0-spec-final-pdf}.
\bibitem{cacoilo2026} N. Cacoilo et al., ``130-nm CMOS-Integrated Superparamagnetic Tunnel Junction-Based p-bit,'' \emph{IEEE Electron Device Letters} 47 (2026), DOI: 10.1109/LED.2026.3696800.
\bibitem{daniel2024} J. Daniel et al., ``Experimental demonstration of an on-chip p-bit core based on stochastic magnetic tunnel junctions and 2D MoS2 transistors,'' \emph{Nature Communications} 15, 4098 (2024), DOI: 10.1038/s41467-024-48152-0.
\end{thebibliography}
\end{document}
